% Calibración del fotómetro micro:bit — del número crudo al significado.
%
% Diagrama propio en TikZ (sin licencias de terceros, vectorial, editable). Se
% tiñe solo con el color de la línea (milcGradeDark). Mismo lenguaje visual que
% la escala de magnitud del módulo 1: una barra donde el sensor entrega un
% número de 0 a 255 y los UMBRALES lo traducen a «oscuro / intermedio / con luz».
% Calibrar = decidir dónde van esas dos líneas.
%
% Se incrusta vía \guiaDiagrama desde el bloque `recursos:` del YAML.
\begin{tikzpicture}[scale=1,font=\sffamily]

  % x = valor/255 * 10.5  →  la barra mide 10.5 cm de ancho
  \def\bx#1{#1/255*10.5}

  % ── Zonas de la barra (calibradas: umbrales en 30 y 90) ──────────────
  \fill[milcGradeDark]      (0,0) rectangle ({\bx{30}},0.75);
  \fill[milcGradeDark!45]   ({\bx{30}},0) rectangle ({\bx{90}},0.75);
  \fill[milcMostaza]        ({\bx{90}},0) rectangle ({\bx{255}},0.75);
  \draw[milcNegro!55,line width=.5pt] (0,0) rectangle ({\bx{255}},0.75);

  % ── Etiquetas de las zonas ───────────────────────────────────────────
  \node[font=\scriptsize\bfseries,white] at ({\bx{15}},0.375) {oscuro};
  \node[font=\scriptsize\bfseries,milcNegro] at ({\bx{60}},0.375) {intermedio};
  \node[font=\scriptsize\bfseries,milcNegro] at ({\bx{172}},0.375) {con luz};

  % ── Escala 0..255 ────────────────────────────────────────────────────
  \foreach \v in {0,30,90,150,210,255} {
    \draw[milcNegro!60] ({\bx{\v}},0) -- ({\bx{\v}},-0.12);
    \node[font=\scriptsize,milcNegro!70,below] at ({\bx{\v}},-0.12) {\v};
  }
  \node[font=\scriptsize\bfseries,milcNegro!70,right] at (10.7,-0.32) {input.lightLevel()};

  % ── Las dos líneas de umbral (lo que el estudiante calibra) ──────────
  \draw[dashed,milcTurquesa,line width=.9pt] ({\bx{30}},-0.35) -- ({\bx{30}},1.55);
  \draw[dashed,milcTurquesa,line width=.9pt] ({\bx{90}},-0.35) -- ({\bx{90}},1.55);
  \node[font=\scriptsize\bfseries,milcTurquesa,align=center] at ({\bx{30}},1.8) {umbral 1\\(30)};
  \node[font=\scriptsize\bfseries,milcTurquesa,align=center] at ({\bx{90}},1.8) {umbral 2\\(90)};

  % ── Una medida real: el ojo dijo «oscuro», el instrumento dice 42 ────
  \draw[->,milcNegro,line width=1pt] ({\bx{42}},2.5) -- ({\bx{42}},0.85);
  \fill[milcNegro] ({\bx{42}},0.85) circle (2pt);
  \node[font=\scriptsize\bfseries,milcNegro,align=center] at ({\bx{42}},2.85)
    {tu medida: \textbf{42}\\(el ojo dijo «oscuro»)};

  % ── Moraleja ─────────────────────────────────────────────────────────
  \node[font=\scriptsize\itshape,milcNegro!75,align=center,text width=10.6cm]
    at (5.25,-1.15)
    {El sensor solo entrega un número de 0 a 255; \textbf{el umbral lo traduce a
     significado}. Calibrar es decidir dónde van esas dos líneas ---y ahí una
     medida de 42 deja de ser «me pareció oscuro» y pasa a ser un dato que otro
     puede comprobar.};

\end{tikzpicture}
